% 核心观点：性能差分准确连接部署动作后果与任务表现，动作MSE只能在附加条件下提供界。
% 叙述逻辑：在完整历史上望远镜求和 -> 覆盖与敏感性给出动作风险界 ->
% 未覆盖情形与经验泛化边界 -> 完整动作历史的结构性失配 -> 动作风险与成功率排序反例。
\section{Task Performance and Imitation Risk}
\label{app:actionproofs}
% \subsection{性能差分恒等式}
\subsection{Performance-Difference Identity}
% 完整历史 $x_t=(s_{0:t},\mathcal H_t)$ 在固定受控过程下构成分析用状态。
The complete history $x_t=(s_{0:t},\mathcal H_t)$ forms an analysis state under the fixed controlled process.
% 运行信息由A节的运行记录；共同过程要求计算与接收转移相同，不只是记录字段相同。
Runtime information is included through $\mathcal R_{0:t}$ as defined in Section~\ref{sec:task}.
The common-process assumption requires the same computation and reception transition law, not merely the same record format.
Internal plans can be reconstructed from the respective causal histories;
% 不将两个策略各自的缓存直接视为同一个物理状态。
their separate caches are not identified with a common physical state.

% 评价结束时的任务失败指示量定义为 $V_T^0(x_T)=1-\Phi(s_{0:T})$，
At the end of evaluation, define the task-failure indicator as $V_T^0(x_T)=1-\Phi(s_{0:T})$.
% 任务成功时取 $0$，失败时取 $1$，其期望为任务失败概率。
It equals $0$ for success and $1$ for failure, and its expectation is the task-failure probability.
% 在 $0\le t<T$ 时，条件期望的递推关系为
For $0\le t<T$, the conditional expectations satisfy the recursion
\[
 Q_t^0(x,a)=
 \E[V_{t+1}^0(x_{t+1})\mid x_t=x,\ u_t=a].
\]
% 此处 $u_t=a$ 表示主动执行动作 $a$，再按共同的状态转移规律推进系统，
Here, $u_t=a$ means actively executing action $a$ and advancing the system under the common transition law,
% 而非仅筛选记录中恰好执行了 $a$ 的样本。
rather than filtering recorded samples in which action $a$ happened to occur.
% 对策略 $\pi$ 的实际轨迹取期望，得到
Taking expectations over trajectories generated by $\pi$ gives
\begin{align}
 \E_\pi\sum_{t=0}^{T-1}A_t^0(x_t,\pi(x_t))
 &=\E_\pi\sum_{t=0}^{T-1}
   [V_{t+1}^0(x_{t+1})-V_t^0(x_t)]\nonumber\\
 &=\E_\pi[V_T^0(x_T)-V_0^0(x_0)]\nonumber\\
 &=F(\pi)-F(\pi_0).
 \label{eq:pdproof}
\end{align}
% 最后一步使用相同初始分布，以及 $V_0^0$ 表示从初始情形起遵循参考策略的失败概率。
The last step uses the common initial distribution and the definition of $V_0^0$ as the failure probability under the reference policy from the initial situation.
% 按时间混合分布改写即得式\eqref{eq:pd}。
Expressing the sum through the time-mixture distribution gives~\eqref{eq:pd}.

% \subsection{动作风险界及其适用条件}
\subsection{Action-Risk Bound and Its Conditions}
% 正文的事件概率覆盖条件与以下测度形式等价：
The event-probability coverage condition in the main text is equivalent to
\[
 \bar\mu_\pi\ll\nu,\qquad
 \frac{\mathrm d\bar\mu_\pi}{\mathrm d\nu}\le C
 \quad \nu\text{-almost everywhere}.
\]
% 其中，$\ll$ 表示绝对连续，$\mathrm d\bar\mu_\pi/\mathrm d\nu$ 为Radon--Nikodym导数，即两分布的密度比。
Here, $\ll$ denotes absolute continuity, and $\mathrm d\bar\mu_\pi/\mathrm d\nu$ is the Radon--Nikodym derivative, or density ratio, of the two distributions.
% 若对所有可测事件均有 $\bar\mu_\pi(\mathcal B)\le C\nu(\mathcal B)$，则 $\nu(\mathcal B)=0$ 蕴含 $\bar\mu_\pi(\mathcal B)=0$，
% 因而上述导数存在；若它在某个正 $\nu$ 概率的集合上超过 $C$，在该集合上积分就会违反事件概率不等式。
If $\bar\mu_\pi(\mathcal B)\le C\nu(\mathcal B)$ for all measurable events, then $\nu(\mathcal B)=0$ implies $\bar\mu_\pi(\mathcal B)=0$, so this derivative exists.
If it exceeded $C$ on a set of positive $\nu$ probability, integrating over that set would violate the event-probability inequality.
% 反之，密度比上界在任意可测事件上积分即得正文的覆盖条件。
Conversely, integrating the density-ratio bound over any measurable event gives the coverage condition in the main text.
% 因此，该条件也约束任意非负可测量的期望；证明先对非负简单函数成立，再由单调收敛推广。
The same condition bounds expectations of nonnegative measurable functions: this follows first for nonnegative simple functions and then by monotone convergence.

% 在正文给定的覆盖与敏感性条件下，记
Under the coverage and sensitivity conditions in the main text, let
% $e(x)=\norm{\pi(x)-\pi_{\mathrm E}(x)}$，则
$e(x)=\norm{\pi(x)-\pi_{\mathrm E}(x)}$. Then
\begin{align}
 F(\pi)-F(\pi_{\mathrm E})
 &\le TL\,\E_{\bar\mu_\pi}e(x)\nonumber\\
 &\le TL\sqrt{\E_{\bar\mu_\pi}e(x)^2}
 \le TL\sqrt{C R_\nu(\pi)} .
 \label{eq:riskproof}
\end{align}
% 第二步使用Cauchy--Schwarz不等式，第三步将上述覆盖条件用于非负量 $e(x)^2$，得到 $\E_{\bar\mu_\pi}e(x)^2\le C\E_\nu e(x)^2$。
The second step uses the Cauchy--Schwarz inequality. The third applies coverage to the nonnegative quantity $e(x)^2$, giving $\E_{\bar\mu_\pi}e(x)^2\le C\E_\nu e(x)^2$.
% 若只能在一个覆盖部分建立密度比界，
If a density-ratio bound holds only on a covered component,
% 例如 $\bar\mu_\pi=(1-\kappa)\nu_{\mathrm{cov}}+\kappa\nu_{\mathrm{out}}$，
for example, $\bar\mu_\pi=(1-\kappa)\nu_{\mathrm{cov}}+\kappa\nu_{\mathrm{out}}$,
% 且 $\mathrm d\nu_{\mathrm{cov}}/\mathrm d\nu\le C$，则由 $A_t^{\mathrm E}\le1$ 得
with $\mathrm d\nu_{\mathrm{cov}}/\mathrm d\nu\le C$, then $A_t^{\mathrm E}\le1$ yields
\[
 F(\pi)-F(\pi_{\mathrm E})
 \le T(1-\kappa)L\sqrt{C R_\nu(\pi)}+T\kappa .
\]
% 其中，$\kappa$ 为未覆盖混合分量的权重。
Here, $\kappa$ is the weight of the uncovered mixture component.
% 该式仍只是上界，不能把 $T\kappa$ 解释为失败率下界；
This remains an upper bound; $T\kappa$ cannot be interpreted as a lower bound on failure probability.
% 失败概率的上界还可与 $1$ 取较小值，因为 $F(\pi)\le1$。
The upper bound may also be capped at $1$, since $F(\pi)\le1$.

% 从经验训练损失到总体风险还需独立的统计论证。
Relating empirical training loss to population risk requires a separate statistical argument.
% 如果在与 $\nu$ 一致的采样及加权协议下，
Suppose that, under a sampling and weighting protocol consistent with $\nu$,
% 已经得到概率至少 $1-\alpha$ 成立的保证
the following guarantee holds with probability at least $1-\alpha$:
$R_\nu(\widehat\pi)\le\widehat R(\widehat\pi)+\epsilon_{\mathrm{gen}}$,
% 则同一事件上
Then, on the same event,
\[
 F(\widehat\pi)\le F(\pi_{\mathrm E})
 +TL\sqrt{C[\widehat R(\widehat\pi)+\epsilon_{\mathrm{gen}}]}.
\]
% 其中，$\widehat R$ 为采用相同动作尺度的经验平方风险，
where $\widehat R$ is the empirical squared risk using the same action scale,
% $\epsilon_{\mathrm{gen}}$ 为对应策略选择过程的有效泛化余项。
and $\epsilon_{\mathrm{gen}}$ is a valid generalization remainder for the corresponding policy-selection procedure.
% 本文不从固定数据集本身推出这一保证；
We do not infer this guarantee from the fixed dataset alone;
% 动作块的混合损失也不能不加转换地代入 $\widehat R$。
nor can a composite action-chunk loss be substituted for $\widehat R$ without establishing the required correspondence.

% \subsection{完整动作历史上的覆盖限制}
\subsection{Coverage Limitations for Complete Action Histories}
\label{app:historycoverage}
% 取 $T\ge2$，确定性专家在初始时刻恒执行 $u_0=0$，
Let $T\ge2$. A deterministic expert always executes $u_0=0$ initially,
% 学习器恒执行一个允许的非零动作 $u_0=\varepsilon$。
whereas the learner always executes an admissible nonzero action $u_0=\varepsilon$.
% 由于 $x_1$ 包含已执行动作，事件
Because $x_1$ contains the executed action, the event
% $\mathcal B=\{t=1,\ u_0=\varepsilon\}$ 满足
$\mathcal B=\{t=1,\ u_0=\varepsilon\}$ satisfies
\[
 \bar\mu_{\pi_{\mathrm E}}(\mathcal B)=0,
 \qquad \bar\mu_\pi(\mathcal B)=1/T>0.
\]
% 故当 $\nu=\bar\mu_{\pi_{\mathrm E}}$ 时，$\bar\mu_\pi\not\ll\nu$，
Thus, for $\nu=\bar\mu_{\pi_{\mathrm E}}$, we have $\bar\mu_\pi\not\ll\nu$,
% 无论 $\varepsilon$ 多小，也无论专家示范增加多少条。
regardless of how small $\varepsilon$ is or how many expert demonstrations are collected.
% 这一反例甚至允许动作不影响下一步物理状态；失配来自记录的动作历史本身。
This counterexample even allows actions to have no effect on the next physical state; the mismatch arises from the recorded action history itself.

% 选取更广的参考分布只可能改变覆盖条件，不能同时提供该分布上的专家动作标签或泛化保证。
Choosing a broader reference distribution can change coverage, but does not simultaneously provide expert action labels or generalization guarantees on that distribution.
% 若改在较小状态上建立覆盖界，还需证明该状态足以表达任务记忆、受控转移和动作块执行规则，
Replacing the analysis state by a reduced state would additionally require proving that it retains the task memory, controlled transitions, and action-chunk execution rules,
% 并据此重新定义沿参考策略执行的失败概率；以下仅在表示空间迁移动作风险，性能差分仍在完整历史上进行。
and redefining the failure probability under reference-policy continuation accordingly. The representation-based risk transfer below instead retains complete histories for the performance-difference identity.

% \subsection{表示空间中的动作风险迁移}
\subsection{Action-Risk Transfer Through a Representation}
\label{app:representationcoverage}
% 该节是补充风险迁移分析，不是C节任务后果表示损失命题的推导前提。
This supplementary risk-transfer analysis is not a premise of the task-consequence representation bound in Section~\ref{sec:representationanalysis}.
% 表示覆盖与信息损失的权衡也见域适应文献，但以下结论由本文的条件独立推导。
Related coverage--information tradeoffs appear in domain adaptation~\cite{johansson2019support,zhao2019invariant}; the following result is derived under the present conditions.
% 固定策略、专家和同一个可测因果编码器，时间索引保留在表示中。
Fix the policy $\pi$, expert $\pi_{\mathrm E}$, and a common measurable causal encoder $Z=\phi(t,\mathcal H_t)$ retaining the time index.
% mu为完整历史部署分布，mu_Z和nu_Z分别为Z的诱导分布。
Let $\mu=\bar\mu_\pi$, and let $\mu_Z$ and $\nu_Z$ denote the induced distributions of $Z$.
% 假设表示空间为标准Borel空间，使以下条件期望具有关于z的可测版本；动作平方误差在mu和nu下可积。
Assume the representation space is standard Borel so that the conditional expectations below admit measurable versions as functions of $z$, and that $e(x)^2=\norm{\pi(x)-\pi_{\mathrm E}(x)}^2$ is integrable under both $\mu$ and $\nu$.
% m_mu和m_nu分别是相同表示下的部署与参考条件平均平方误差，D_phi度量两者的部署加权差异。
Define the conditional mean squared errors and their deployment-weighted discrepancy by
\begin{align*}
 m_\mu(z)&=\E_\mu[e(x)^2\mid Z=z],\\
 m_\nu(z)&=\E_\nu[e(x)^2\mid Z=z],\\
 D_\phi&=\E_{z\sim\mu_Z}|m_\mu(z)-m_\nu(z)|.
\end{align*}
% 假设对每个可测事件有表示覆盖；它使m_nu的参考零概率集合在部署下也为零概率，因而D_phi不依赖版本选择。
Suppose $\mu_Z(\mathcal B)\le C_\phi\nu_Z(\mathcal B)$ for every measurable event, with $1\le C_\phi<\infty$.
This condition ensures that reference null sets are also deployment null sets, so $D_\phi$ is independent of the versions chosen for the conditional means.
% 由条件期望及非负函数的覆盖界，
Conditioning on $Z$ and applying coverage to the nonnegative function $m_\nu$ give
\begin{align}
 \E_\mu e(x)^2
 &=\E_{\mu_Z}m_\mu(Z)\nonumber\\
 &\le \E_{\mu_Z}m_\nu(Z)+D_\phi\nonumber\\
 &\le C_\phi\E_{\nu_Z}m_\nu(Z)+D_\phi\nonumber\\
 &=C_\phi R_\nu(\pi)+D_\phi.
 \label{eq:representationrisktransfer}
\end{align}
% 有限覆盖常数和两边平方可积也保证D_phi有限。再应用敏感性与Cauchy--Schwarz得正文结果。
These assumptions also make $D_\phi$ finite and establish~\eqref{eq:representationrisktransfer-main}.
Under the sensitivity condition in Section~\ref{sec:actionerrorbound}, the performance-difference identity gives
$F(\pi)-F(\pi_{\mathrm E})\le TL\E_\mu e(x)$.
Applying Cauchy--Schwarz and~\eqref{eq:representationrisktransfer} yields the main-text bound~\eqref{eq:representationcoveragebound}.
% 这里不需要完整历史覆盖，也不需要表示具有Markov性。
This does not require complete-history coverage or a Markov assumption on $Z$.
% 若存在相同的可测函数使策略和专家在两分布下均只依赖Z，则e平方是Z的共同函数，D_phi为零。
If common measurable functions express both $\pi(x)$ and $\pi_{\mathrm E}(x)$ solely in terms of $Z$ under both distributions, then $e(x)^2$ is a common function of $Z$ and $D_\phi=0$.
% 一般情形下D_phi并非单一歧义指标，零值也不推出表示充分；它只是该固定策略的条件平均误差差异。
In general, $D_\phi$ measures a conditional mean-error difference for this fixed policy, not all decision ambiguity; $D_\phi=0$ alone does not establish representation sufficiency.
% D_phi与C节基于任务后果的表示损失不同，也涉及离线示范无法单独提供的部署专家动作。
It differs from the task-consequence representation loss in Section~\ref{sec:representationanalysis} and involves expert actions on deployment histories that fixed demonstrations alone do not supply.
% 该界是固定对象的条件性界，而不是数据依赖训练过程的统一泛化保证。
The result is a conditional bound for fixed objects, not a uniform generalization guarantee for data-dependent training.

% 可测压缩将事件合并，因此完整空间的覆盖常数也适用于表示空间；反方向通常不成立。
Measurable compression merges events: any complete-space coverage constant also bounds representation events by taking their preimages, but the converse need not hold.
% 可测可逆变换在两方向均保留事件，故最佳覆盖常数不变。保留RGB并附加其确定性语义图也有可测逆投影。
A measurable invertible transformation with measurable inverse preserves the best coverage constant in both directions.
Appending a deterministic semantic map to retained RGB has the same property relative to RGB, since projection recovers the input.
% 仅在有损表示时覆盖可改善，而式中的D_phi保留条件误差差异的代价；历史增加也不保证覆盖常数更小。
Coverage can improve through a lossy representation, with $D_\phi$ retaining the cost of conditional error differences; adding history need not reduce the coverage constant.

% \subsection{更小动作风险但更差任务表现的反例}
\subsection{A Counterexample: Lower Action Risk but Worse Task Performance}
% 考虑 $T=1$、两个等概率且可观测的初始状态，
Consider $T=1$ with two equally likely, observable initial states,
% 动作集合为 $[0,0.2]$，专家在两种状态均选择 $0$。
action set $[0,0.2]$, and an expert that chooses $0$ in both states.
% 第一种状态的失败概率为 $\min(1,100a^2)$，
The failure probability in the first state is $\min(1,100a^2)$,
% 第二种状态始终成功。
while the second state always leads to success.
% 终局图像可直接显示成功或失败，因此满足视觉可评估性。
The final image directly reveals success or failure, satisfying visual assessability.

% 策略 $\pi_1$ 分别输出 $(0.1,0)$，策略 $\pi_2$ 分别输出 $(0,0.2)$。
Policies $\pi_1$ and $\pi_2$ output $(0.1,0)$ and $(0,0.2)$ across the two states, respectively.
% 在共同初始分布 $\nu$ 上，
Under the common initial distribution $\nu$,
\begin{align}
 R_\nu(\pi_1)&=0.005<0.02=R_\nu(\pi_2),\nonumber\\
 J(\pi_1)&=0.5<1=J(\pi_2).
 \label{eq:counterexample}
\end{align}
% 本例中覆盖常数 $C=1$，敏感性界可取 $L=10$，
The coverage constant is $C=1$, and the sensitivity bound holds with $L=10$,
% 因为 $\min(1,100a^2)\le10a$。
since $\min(1,100a^2)\le10a$.
% 因而即使正文的界成立，较小动作风险也不蕴含策略成功率排序。
Thus, even when the bound in the main text holds, smaller action risk does not imply higher task success.
% 差异来自误差发生的位置与任务后果，而不是界的代数错误。
The difference comes from where errors occur and their task consequences, not from an algebraic error in the bound.
